\documentclass[hidelinks,onefignum,onetabnum,final]{siamart220329}  % for arxiv
%\documentclass[review,hidelinks,onefignum,onetabnum,final]{siamart220329}  % for submission

\usepackage{amsfonts,amssymb}
\usepackage{graphicx}
\DeclareGraphicsExtensions{.eps,.pdf,.png,.jpg}

% Used for creating new theorem and remark environments
\newsiamremark{remark}{Remark}
%definition already defined \newsiamremark{definition}{Definition}
\newsiamremark{hypothesis}{Hypothesis}
\crefname{hypothesis}{Hypothesis}{Hypotheses}
\newsiamthm{claim}{Claim}
\newsiamremark{example}{Example}

\usepackage{bm}

%\usepackage{amsopn}
%\usepackage{bm,bbm,empheq,verbatim,fancyvrb}
%\usepackage{booktabs,multirow,xspace}

\usepackage{tikz}
\usetikzlibrary{decorations.pathreplacing}
\usetikzlibrary{graphs,quotes}

\usepackage[kw]{pseudo}
\pseudoset{%
left-margin=6mm,%
topsep=5mm,%
label=\footnotesize\arabic*,%
idfont=\texttt,%
ctfont=\textsl,%
ct-left=\qquad\qquad(,%
ct-right=),%
}

\newcommand{\eps}{\epsilon}
\newcommand{\RR}{\mathbb{R}}

\newcommand{\grad}{\nabla}
\newcommand{\Div}{\nabla\cdot}

\newcommand{\bn}{\mathbf{n}}
\newcommand{\bu}{\mathbf{u}}
\newcommand{\bw}{\mathbf{w}}

\newcommand{\bU}{\mathbf{U}}

\newcommand{\bPhi}{\bm{\Phi}}
\newcommand{\bzero}{\bm{0}}

\newcommand{\dx}{\,dx}
\newcommand{\bds}{\bm{ds}}

\newcommand{\cK}{\mathcal{K}}
\newcommand{\cU}{\mathcal{U}}
\newcommand{\cV}{\mathcal{V}}
\newcommand{\cW}{\mathcal{W}}
\newcommand{\cX}{\mathcal{X}}
\newcommand{\cY}{\mathcal{Y}}
\newcommand{\cZ}{\mathcal{Z}}

\newcommand{\pp}{\mathrm{p}}
\newcommand{\qq}{\mathrm{q}}

\newcommand{\ip}[2]{\left<#1,#2\right>}
\newcommand{\Bform}[2]{B\big(\left(#1\right),\left(#2\right)\big)}

\DeclareMathOperator{\diag}{diag}

\newcommand{\XX}{\ding{55}}

\DeclareMathOperator*{\argmin}{arg\,min}

\newfloat{pseudofloat}{t}{xyz}[section]
\floatname{pseudofloat}{Algorithm}

\usepackage{todonotes}
\newcommand{\elb}[1]{\todo[inline, color=blue!20]{EB: #1}}


% Sets running headers as well as PDF title and authors
%\headers{FIXME}{E. Bueler}

\title{Speculation on a fixed-point argument for \\ time-dependent glacier evolution with Stokes dynamics}

\author{Ed Bueler and Claude}

%\date{\today}
%\thanks{Department of Mathematics and Statistics, University of Alaska Fairbanks, USA (\email{elbueler@alaska.edu}).}

\graphicspath{{figs/}}

\begin{document}

\maketitle

\begin{abstract}
This is rank speculation generated by claude when I asked it to do that.  Various links don't resolve because the point into pile.tex.
\end{abstract}

%\begin{keywords}
%FIXME
%\end{keywords}
%\begin{MSCcodes}
%FIXME
%\end{MSCcodes}

FIXME recent references from Claude search:

Daganzo (2005). "A variational formulation of kinematic waves," Transportation Research B 39(1), 49--72.

Lions (1982). "Generalized Solutions of Hamilton-Jacobi Equations". does this treat obstacle problems for first-order equations?

Díaz \& Schiavi (1999). "On a degenerate parabolic/hyperbolic system in glaciology giving rise to a free boundary," Nonlinear Anal. 38, 649--671.  This paper looks at a complementarity problem for the ice bed, with the free boundary being between the frozen part of the bed and the wet part.  The mass continuity problem for the glacier does appear, but has thickness assumed positive everywhere in the domain.  The theorem which is proved is interesting, namely that an implicit time stepping scheme converges to show existence of the continuum time-dependent problem.

\bigskip

Lemma~\ref{lem:skoro} solves complementarity problem~\eqref{eq:pilencp}
when the velocity $(u,v,w)$ is prescribed as data.  In the full glacier
evolution model the velocity is not data---it solves a Stokes sub-model
within the evolving ice domain, so the velocity and surface elevation
are coupled.  This appendix sketches how one might build a fixed-point
existence argument for the coupled problem using the Skorohod formula,
and identifies the key analytical difficulties.  The tone is deliberately
speculative; the goal is to record what would need to be proved.

\subsection{The fixed-point map}

Fix a time horizon $T>0$ and an exponent $r > 2$.  Define the admissible
set of surface elevations
\begin{equation}
\cK = \left\{s \in W^{1,r}(\Omega) : s \ge b \;\text{on}\; \Omega,\;
         s = b \;\text{on}\; \partial\Omega\right\}.
\label{eq:app:admissible}
\end{equation}
Since $\Omega \subset \RR^2$ and $r > 2$, Sobolev's embedding theorem gives
$W^{1,r}(\Omega) \hookrightarrow C(\bar\Omega)$, so each $s\in\cK$ is
continuous.  For $s\in\cK$ define the 3D ice domain
\begin{equation}
\Lambda(s) = \{(x,y,z) : b(x,y) < z < s(x,y)\} \subset \Omega\times\RR;
\label{eq:app:domain}
\end{equation}
continuity of $s$ and $b$ makes this a Lipschitz domain.

Let $\cU$ be a Banach space of surface velocity fields
$\bU = (u,v,w): [0,T]\times\Omega\to\RR^3$, to be specified.  Define
$T = T_2 \circ T_1$, the composition of two steps:
\begin{itemize}
\item $T_1:\cU \to C([0,T];\cK)$: given $\bU \in \cU$, apply
Lemmas~\ref{lem:chars} and~\ref{lem:skoro} to produce the surface
elevation $s = T_1 \bU$.
\item $T_2:C([0,T];\cK) \to \cU$: given $s\in C([0,T];\cK)$, for
each $t\in[0,T]$ solve the non-Newtonian Stokes sub-model
\cite{GreveBlatter2009} on the domain $\Lambda(s(t,\cdot))$,
and return the surface-trace velocity $\bu|_s(t,\cdot)$ extended by
zero to all of $\Omega$.
\end{itemize}
A fixed point $\bU = T\bU$ is a velocity field consistent with the
full coupled Stokes--surface kinematical complementarity problem.

\subsection{The Skorohod step \texorpdfstring{$T_1$}{T1}}

\paragraph{Input regularity needed.}
Lemma~\ref{lem:chars} requires $u,v \in C^1([0,T]\times\Omega)$ for
the characteristic ODE \eqref{eq:chars} to be classically well-posed.
A weaker sufficient condition is given by the DiPerna--Lions theory
of renormalized flows \cite{DiPernaLions1989}: a unique Lagrangian
flow exists for almost every initial point whenever
$(u,v) \in L^1(0,T;W^{1,1}(\Omega))$ and
$\mathrm{div}(u,v) \in L^1([0,T]\times\Omega)$.
Stokes incompressibility ($\Div\bu=0$ inside $\Lambda$) will be
relevant here.

\paragraph{Spatial regularity of $s$.}
The Jacobian of the characteristic flow map
$\Phi_t:(x_0,y_0)\mapsto(X(t;x_0,y_0),Y(t;x_0,y_0))$ satisfies the
Liouville identity
\begin{equation}
J(t) = \exp\!\left(\int_0^t \mathrm{div}(u,v)\big|_{\Phi_\tau} \, d\tau\right),
\label{eq:app:jacobian}
\end{equation}
so $J(t)>0$ when $\mathrm{div}(u,v) \in L^1$, ensuring invertibility.
When $(u,v)\in L^1(0,T;W^{1,r}(\Omega))$, regularity theory for ODE
flows gives $\Phi_t \in W^{1,r}(\Omega)$, and then the reassembled
surface elevation satisfies
\begin{equation}
s(t,\cdot) \in W^{1,r}(\Omega), \qquad
\|s(t,\cdot)\|_{W^{1,r}} \le
C\!\left(s_0,\, \|\bU\|_{L^1(0,T;\,W^{1,r})},\, \|a\|_{L^1(0,T;\,W^{1,r})}\right),
\label{eq:app:T1bound}
\end{equation}
provided $s_0\in W^{1,r}(\Omega)$ and $c = w+a \in L^1(0,T;W^{1,r}(\Omega))$.

\begin{remark}
The Skorohod correction $\ell(t;x_0,y_0)$ defined by \eqref{eq:skoro}
is a running maximum over $r\in[0,t]$, so its spatial regularity in
$(x_0,y_0)$ requires extra care near times when the maximum is
attained (i.e., where $S = b\circ\Phi$, the glacier margin).
A complete proof of \eqref{eq:app:T1bound} would need to verify,
for instance, that the contact set has measure-zero level sets,
or work with distributional derivatives.
\end{remark}

\subsection{The Stokes step \texorpdfstring{$T_2$}{T2}}

For $s\in\cK\subset W^{1,r}(\Omega)$ with $r>2$, the domain
$\Lambda(s)$ is Lipschitz and the regularized non-Newtonian Stokes
system is well-posed \cite{JouvetRappaz2011}.
Let $\pp = (n+1)/n$ denote the Glen flow exponent; for
$n=3$ one has $\pp=4/3$.  The solution velocity
$\bu\in W_0^{1,\pp}(\Lambda;\RR^3)$ satisfies an a priori bound
\begin{equation}
\|\bu\|_{W^{1,\pp}(\Lambda)} \le C_1,  \label{eq:app:apriori}
\end{equation}
where $C_1$ depends on physical constants and the geometry of
$\Lambda(s)$, in particular on $\|s\|_{W^{1,r}}$; see \cite{Bueler2025}.
A trace inequality then gives
\begin{equation}
\int_{\Gamma_s} |\bu|_s|^{\pp} \, dS \le C_1,  \label{eq:app:trace}
\end{equation}
and extended by zero, the surface trace satisfies
$\bu|_s(t,\cdot) \in L^{\pp}(\Omega;\RR^3)$ for each $t$.

\paragraph{What $T_2$ delivers.}
$T_2$ returns $\bu|_s \in L^{\pp}(\Omega;\RR^3)$, i.e.~velocity
components in $L^{4/3}(\Omega)$ for $n=3$.
Crucially, $T_2$ does \emph{not} directly deliver $W^{1,q}(\Omega)$
regularity for any $q > 0$.

\subsection{The regularity gap}

For the cycle
$\cU \xrightarrow{T_1} C([0,T];\cK) \xrightarrow{T_2} \cU$
to close, $\cU$ must satisfy both the input requirement of $T_1$
and the output property of $T_2$.  Since $T_1$ requires at minimum
$(u,v)\in L^1(0,T;W^{1,1}(\Omega))$ (for DiPerna--Lions) but
$T_2$ delivers only $\bu|_s \in L^1(0,T;L^{\pp}(\Omega))$
with $\pp=4/3$, the cycle does not close:
\[
W^{1,r}\text{ velocity} \;\xrightarrow{T_1}\; W^{1,r}\text{ surface}
\;\xrightarrow{T_2}\; L^{\pp}\text{ velocity}
\;\not\hookrightarrow\; W^{1,1}\text{ velocity}.
\]

There is also an exponent tension.  The surface elevation space requires
$r > 2$ (for $\Lambda(s)$ to be a Lipschitz domain via Sobolev
embedding), while the Stokes law has exponent $\pp = 4/3 < 2$.
Even if $\bu|_s \in W^{1,\pp}(\Omega)$ could be established, one has
$W^{1,\pp} \not\subset W^{1,r}$ for $r > \pp$, so the input
requirement of $T_1$ would still not be met by the output of $T_2$.

\subsection{Hypotheses that would close the argument}

Gaining spatial regularity of $\bu|_s$ beyond $L^{\pp}$ would
require higher interior regularity of the Stokes velocity:
$W^{2,\pp}$ regularity of $\bu$ inside $\Lambda(s)$ would give
(by a trace theorem) $\bu|_s \in W^{1,\pp}(\Omega)$.
Such $W^{2,\pp}$ regularity requires $\partial\Lambda(s)$ to be
$C^{1,\alpha}$ \cite{JouvetRappaz2011}, i.e.~$s\in C^{1,\alpha}(\Omega)$,
which in turn follows from $s\in W^{2,r}(\Omega)$ via Sobolev.
So higher surface elevation regularity would bootstrap Stokes
regularity, but the argument is circular without an independent
source of that higher regularity.

\begin{hypothesis}[Stokes surface trace regularity]
\label{hyp:stokestraceW1}
There exist exponents $r>2$ and $q>1$ such that for every
$s\in\cK \subset W^{1,r}(\Omega)$ the Stokes velocity surface trace
satisfies $\bu|_s \in W^{1,q}(\Omega;\RR^3)$, with
\begin{equation}
\|\bu|_s\|_{W^{1,q}(\Omega)} \le C\bigl(\|s\|_{W^{1,r}(\Omega)}\bigr).
\label{eq:app:hyp1}
\end{equation}
\end{hypothesis}

\noindent
For \emph{Newtonian} ice ($n=1$, $\pp=2$) on a $C^{1,1}$ domain,
classical $H^2$ regularity of the linear Stokes system gives
$\bu|_s \in H^1(\Omega)$, so Hypothesis~\ref{hyp:stokestraceW1} holds
with $q=2$ in that case.
For shear-thinning ice ($\pp<2$) on a domain that is merely Lipschitz,
this appears to be open.
A related domain-dependence statement for the Stokes map is
Conjecture~1 in \cite{Bueler2025}; that conjecture concerns
Lipschitz continuity rather than regularity, but both ultimately
rest on the same missing theory of how the Stokes solution varies
with the domain.

If Hypothesis~\ref{hyp:stokestraceW1} holds with $q\ge 1$, and if
the surface trace $(u,v) = (\bu|_s)_{1,2}$ satisfies
$\mathrm{div}(u,v) \in L^1$ (which follows from Stokes
incompressibility via the relation $w|_s = -\int_b^s \Div(u,v)\,dz$,
up to a boundary term at the margin), then the DiPerna--Lions
theory \cite{DiPernaLions1989} applies to the characteristic ODE
\eqref{eq:chars}, closing the regularity loop.

\begin{hypothesis}[Lipschitz continuity of Stokes trace in $s$]
\label{hyp:lipschitz}
Under the conditions of Hypothesis~\ref{hyp:stokestraceW1}, for every
$R>0$ there exists $C(R)>0$ such that for $s_1,s_2\in B_R\cap\cK$,
\begin{equation}
\|\bu|_{s_1} - \bu|_{s_2}\|_{W^{1,q}(\Omega)} \le
C(R)\,\|s_1-s_2\|_{W^{1,r}(\Omega)}.  \label{eq:app:hyp2}
\end{equation}
\end{hypothesis}

\noindent
Hypothesis~\ref{hyp:lipschitz} with $W^{1,q}$ replaced by $L^{r'}$
(where $1/r+1/r'=1$) is exactly Conjecture~1 in \cite{Bueler2025},
which is also the key ingredient for well-posedness of the implicit
time-step variational inequality in that reference.

\paragraph{Invariant set.}
For either the Schauder or Banach fixed-point theorem one needs a
bounded closed convex set $K\subset\cU$ with $T(K)\subseteq K$.
The a priori Stokes bound \eqref{eq:app:apriori} gives
$C_1=C_1(R_s)$ with $R_s = \|s\|_{W^{1,r}}$, and
\eqref{eq:app:T1bound} controls $R_s$ in terms of
$R_\cU=\|\bU\|_\cU$.
A self-consistent radius $R>0$ therefore exists if
\begin{equation}
C_1\bigl(C_{\mathrm{Sko}}(R,s_0,a)\bigr) \le R,
\label{eq:app:selfconsistent}
\end{equation}
where $C_{\mathrm{Sko}}$ is the bound from \eqref{eq:app:T1bound}.
For small time horizon $T$, or small forcing $a$, such a radius
can always be found.

\paragraph{Compactness vs.\ contraction.}
Under Hypotheses~\ref{hyp:stokestraceW1} and~\ref{hyp:lipschitz},
$T$ is continuous on bounded sets.
A \emph{Schauder fixed point} additionally requires $T$ to be compact;
compactness of $T_2$ follows from the uniform $W^{1,\pp}(\Lambda)$
bound on $\bu$ together with the compact Sobolev embedding
$W^{1,\pp}\hookrightarrow\hookrightarrow L^m$ for $m<\pp^*$,
but propagating that compactness through $T_1$ requires further
argument.
A \emph{Banach contraction mapping} requires
$\|T\bU_1-T\bU_2\|_\cU \le L\|\bU_1-\bU_2\|_\cU$ with $L<1$;
this would follow from Hypothesis~\ref{hyp:lipschitz} combined with
Gronwall-type ODE estimates for the Lipschitz continuity of $T_1$,
provided $T$ (or the time-step $\Delta t$) is small enough.

\subsection{Relation to the fully implicit approach in \cite{Bueler2025}}  \label{app:implicit}

The framework in \cite{Bueler2025} avoids the above regularity cycle
by formulating an \emph{implicit} time-step variational inequality (VI):
find the new surface elevation $s\in\cK$ such that the Stokes velocity
at $s$ and $s$ itself simultaneously satisfy the VI.
Well-posedness of this implicit step is established in \cite{Bueler2025}
(conditionally on two conjectures analogous to
Hypotheses~\ref{hyp:stokestraceW1}--\ref{hyp:lipschitz}) using the
theory of coercive monotone operators \cite{KinderlehrerStampacchia1980},
without any Picard iteration between Skorohod and Stokes steps.

The Skorohod-based iteration proposed here---alternating between
the Skorohod formula (with given velocity) and the Stokes solve
(with given surface)---is an alternative route to the same coupled
solution.
Convergence of the iteration follows when $T$ is a contraction; under
uniqueness of the implicit VI solution
(which follows from strict monotonicity under Conjecture~2 of
\cite{Bueler2025}), the fixed-point limit and the implicit VI solution
then coincide.

In summary, the central open problem for either approach is
Hypothesis~\ref{hyp:stokestraceW1}: that the Stokes velocity surface
trace on a Lipschitz glacier domain has $W^{1,q}$ regularity as a
function of the map-plane coordinates, with a bound that is Lipschitz
in the surface elevation.

\section{Does Rodrigues use the Skor\-okhod formula?}

Having read \cite{Rodrigues2004} in full, here is an honest comparison
with the approach in pile.tex.

Rodrigues is the closest prior work, but does not use the Skor\-okhod
formula.  He treats the obstacle problem for a first-order operator
$Hu = b\cdot\nabla u + b_0 u$
with constraint $u \ge \psi$, formulated as the complementarity problem
(his eq.~(3.1)):
\[
u \ge \psi, \qquad Hu - f \ge 0, \qquad (Hu-f)(u-\psi) = 0
\quad \text{a.e.\ in } Q.
\]
This is structurally the Skor\-okhod condition distributed over the
domain, but Rodrigues solves it via penalization and variational
inequality methods (the Mignot--Puel 1976 framework), not via the
running-supremum formula.

The reason is that Rodrigues requires the coercivity condition
$b_0 - \tfrac{1}{2}\nabla\cdot b \ge \beta > 0$
for $L^2$ well-posedness.  That zeroth-order damping breaks the clean
characteristic reduction: along a characteristic $\dot X = b(X)$ the
equation acquires an exponential weight, so it is no longer pure
transport.

For pure transport $\partial_t u + b\cdot\nabla u = f$ with $b_0=0$
and obstacle $u \ge 0$, along each characteristic $\dot X = b(X)$
the problem does collapse to the one-dimensional Skor\-okhod problem,
and the solution is $u(X(t),t) = x(t) + \sup_{s\le t} x^-(s)$ fiber
by fiber.  This appears to be a folk theorem: the argument is immediate
from characteristics together with the Skor\-okhod formula, but no
paper states it explicitly by that name.

The foundational works cited in \cite{Rodrigues2004} are Mignot--Puel
(1976) and Bensoussan--Lions (1973); for conservation law obstacle
problems, L\'evi (M2AN, 2001).  None invoke Skor\-okhod by name.


\bibliographystyle{siamplain}
\bibliography{pile}

\end{document}
